\documentclass[10pt,a4paper]{amsart}
\usepackage{amsmath,amssymb,amsthm,graphicx}
\usepackage{verbatim}
\newenvironment{reference}{\comment}{\endcomment}
%\newenvironment{reference}{}{}
\newenvironment{slogan}{\comment}{\endcomment}
\newenvironment{history}{\comment}{\endcomment}

% For commutative diagrams we use Xy-pic
\usepackage[all]{xy}

% We use 2cell for 2-commutative diagrams.
\xyoption{2cell}
% Theorem environments.
%
\theoremstyle{plain}
\newtheorem{theorem}[subsection]{Theorem}
\newtheorem{proposition}[subsection]{Proposition}
\newtheorem{lemma}[subsection]{Lemma}

\theoremstyle{definition}
\newtheorem{definition}[subsection]{Definition}
\newtheorem{example}[subsection]{Example}
\newtheorem{exercise}[subsection]{Exercise}
\newtheorem{situation}[subsection]{Situation}

\theoremstyle{remark}
\newtheorem{remark}[subsection]{Remark}
\newtheorem{remarks}[subsection]{Remarks}

\numberwithin{equation}{subsection}

% Macros
%
\def\lim{\mathop{\mathrm{lim}}\nolimits}
\def\colim{\mathop{\mathrm{colim}}\nolimits}
\def\Spec{\mathop{\mathrm{Spec}}}
\def\Hom{\mathop{\mathrm{Hom}}\nolimits}
\def\Ext{\mathop{\mathrm{Ext}}\nolimits}
\def\SheafHom{\mathop{\mathcal{H}\!\mathit{om}}\nolimits}
\def\SheafExt{\mathop{\mathcal{E}\!\mathit{xt}}\nolimits}
\def\Sch{\mathit{Sch}}
\def\Mor{\mathop{\mathrm{Mor}}\nolimits}
\def\Ob{\mathop{\mathrm{Ob}}\nolimits}
\def\Sh{\mathop{\mathit{Sh}}\nolimits}
\def\NL{\mathop{N\!L}\nolimits}
\def\CH{\mathop{\mathrm{CH}}\nolimits}
\def\proetale{{pro\text{-}\acute{e}tale}}
\def\etale{{\acute{e}tale}}
\def\QCoh{\mathit{QCoh}}
\def\Ker{\mathop{\mathrm{Ker}}}
\def\Im{\mathop{\mathrm{Im}}}
\def\Coker{\mathop{\mathrm{Coker}}}
\def\Coim{\mathop{\mathrm{Coim}}}

% Boxtimes
%
\DeclareMathSymbol{\boxtimes}{\mathbin}{AMSa}{"02}

%
% Macros for moduli stacks/spaces
%
\def\QCohstack{\mathcal{QC}\!\mathit{oh}}
\def\Cohstack{\mathcal{C}\!\mathit{oh}}
\def\Spacesstack{\mathcal{S}\!\mathit{paces}}
\def\Quotfunctor{\mathrm{Quot}}
\def\Hilbfunctor{\mathrm{Hilb}}
\def\Curvesstack{\mathcal{C}\!\mathit{urves}}
\def\Polarizedstack{\mathcal{P}\!\mathit{olarized}}
\def\Complexesstack{\mathcal{C}\!\mathit{omplexes}}
% \Pic is the operator that assigns to X its picard group, usage \Pic(X)
% \Picardstack_{X/B} denotes the Picard stack of X over B
% \Picardfunctor_{X/B} denotes the Picard functor of X over B
\def\Pic{\mathop{\mathrm{Pic}}\nolimits}
\def\Picardstack{\mathcal{P}\!\mathit{ic}}
\def\Picardfunctor{\mathrm{Pic}}
\def\Deformationcategory{\mathcal{D}\!\mathit{ef}}
\usepackage{cleveref}
\crefname{originallemma}{original lemma}{original lemmas}
\Crefname{originallemma}{Original lemma}{Original lemmas}
\makeatletter
\newlabel{chow-lemma-lci-gysin-easy}{{0FF7}{source}{}{sourcepointer-1}{}}
\newlabel{chow-lemma-lci-gysin-easy@cref}{{[originallemma][0][]0FF7}{[1][0][]source}{}{}{}}
\makeatother
\crefname{originalreference}{original reference}{original references}
\Crefname{originalreference}{Original reference}{Original references}
\makeatletter
\newlabel{lemma-composition-correspondences}{{0FG0}{source}{}{sourcepointer-2}{}}
\newlabel{lemma-composition-correspondences@cref}{{[originalreference][0][]0FG0}{[1][0][]source}{}{}{}}
\makeatother
\crefname{originalreference}{original reference}{original references}
\Crefname{originalreference}{Original reference}{Original references}
\makeatletter
\newlabel{lemma-functor-and-cycles}{{0FG5}{source}{}{sourcepointer-3}{}}
\newlabel{lemma-functor-and-cycles@cref}{{[originalreference][0][]0FG5}{[1][0][]source}{}{}{}}
\makeatother
\title{Commutativity of the duality diagram for the transpose correspondence}
\author{The Stacks Project Authors}
\date{stacks-project:weil.tex:local-snapshot-20261012; GNU FDL 1.2 or later}
\begin{document}
\maketitle
\noindent\textit{Editorial note.} All delivered mathematical text is the author's own, copied from the pinned Stacks Project chapter file weil.tex; no statement or proof step was rewritten, reordered or removed.\par\medskip
\noindent\textit{Editorial note.} The counted proof is the author's complete argument: the translation of the diagram into pullbacks via the correspondence of a graph, the two explicit computations, and the identification with the graph of the morphism.\par\medskip
\noindent\textit{Editorial note.} The label lemma-prep-dual names the counted statement, and the proof environment immediately following it in the source is delivered verbatim as its proof.\par\medskip
\noindent\textit{Editorial note.} The counted statement is carried by the author's own lemma environment in the chapter source.\par\medskip
\noindent\textit{Editorial note.} Pre-open gate: python operations/agent-authoring/tools/precard\_gate.py --source weil.tex --label lemma-prep-dual --start 776 --end 827 --local-id 28147 returned RESULT=OK (L28147 free, occupancy FREE, public interval clear). Pre-assembly tag self-check: canonical tag resolved and not present in reference/public-items.jsonl.\par\medskip
\noindent\textit{Editorial note.} The statement and proof are copied from the pinned chapter file; the cited inputs are kept as external pointers with their canonical Stacks tags.\par\medskip
\noindent\textit{Allowed standard prerequisites (editorial).}\begin{itemize}
\item Correspondences of negative degree as cycle groups and composition as pullback. Same chapter, the identification of the correspondence groups from a scheme to the base with cycle groups, and the description of composition with a correspondence as pullback of cycles along it. Turns the commutativity of the diagram into an equality of two pullbacks of the fundamental cycles.
\item Conversion of pushforward and pullback along a graph into the usual operations. Same chapter, the lemma describing pushforward and pullback by the correspondence of a graph in terms of the usual Gysin and pushforward maps. Identifies the two blown up expressions with the graph of the morphism.
\item The Gysin pullback of the diagonal along a graph. Chow Homology, the lemma computing the Gysin pullback of the diagonal along the graph of a morphism. Supplies the last equality identifying the second route around the diagram with the graph of the morphism.
\end{itemize}\medskip
\ifdefined\hypertarget\hypertarget{sourcepointer-1}{}\fi
\noindent\textit{Source reference.} Original lemma 0FF7: sources/stacks-chow/tags.txt, source lines 19277--19277. Prerequisite pointer invoked at proof line(s) 823; it names a standard fact of a sibling chapter that the counted argument uses, and every step of the counted proof built on it is retained verbatim in the delivered body.\par\medskip
\ifdefined\hypertarget\hypertarget{sourcepointer-2}{}\fi
\noindent\textit{Source reference.} Original reference 0FG0: sources/stacks-weil/pinned-original-source.tex, source lines 441--441. Prerequisite pointer invoked at proof line(s) 806; it names a standard fact of the same chapter that the counted argument uses, and every step of the counted proof built on it is retained verbatim in the delivered body.\par\medskip
\ifdefined\hypertarget\hypertarget{sourcepointer-3}{}\fi
\noindent\textit{Source reference.} Original reference 0FG5: sources/stacks-weil/pinned-original-source.tex, source lines 661--661. Prerequisite pointer invoked at proof line(s) 807; it names a standard fact of the same chapter that the counted argument uses, and every step of the counted proof built on it is retained verbatim in the delivered body.\par\medskip
\begin{lemma}
\label{lemma-prep-dual}
Let $f : Y \to X$ be a morphism of smooth projective schemes over $k$.
Assume $X$ and $Y$ equidimensional of dimensions $d$ and $e$.
Denote $a = [\Gamma_f] \in \text{Corr}^0(X, Y)$ and
$a^t = [\Gamma_f^t] \in \text{Corr}^{d - e}(Y, X)$.
Set
$\eta_X = [\Gamma_{X \to X \times X}] \in \text{Corr}^0(X \times X, X)$,
$\eta_Y = [\Gamma_{Y \to Y \times Y}] \in \text{Corr}^0(Y \times Y, Y)$,
$[X] \in \text{Corr}^{-d}(X, \Spec(k))$, and
$[Y] \in \text{Corr}^{-e}(Y, \Spec(k))$. The diagram
$$
\xymatrix{
X \otimes Y \ar[r]_{a \otimes \text{id}} \ar[d]_{\text{id} \otimes a^t} &
Y \otimes Y \ar[r]_{\eta_Y} &
Y \ar[d]^{[Y]} \\
X \otimes X \ar[r]^{\eta_X} &
X \ar[r]^{[X]} &
\Spec(k)
}
$$
is commutative in the category of correspondences.
\end{lemma}

\begin{proof}
Recall that $\text{Corr}^r(W, \Spec(k)) = \CH_{-r}(W)$ for any
smooth projective scheme $W$ over $k$
and given $c \in \text{Corr}^s(W', W)$ the composition
with $c$ agrees with pullback by $c$ as a map
$\CH_{-r}(W) \to \CH_{-r - s}(W')$
(Lemma \ref{lemma-composition-correspondences}).
Finally, we have Lemma \ref{lemma-functor-and-cycles}
which tells us how to convert this into usual
pushforward and pullback of cycles.
We have
$$
(a \otimes \text{id})^* \eta_Y^* [Y] =
(a \otimes \text{id})^* [\Delta_Y] =
(f \times \text{id})_*\Delta_Y = [\Gamma_f]
$$
and the other way around we get
$$
(\text{id} \otimes a^t)^* \eta_X^* [X] =
(\text{id} \otimes a^t)^* [\Delta_X] =
(\text{id} \times f)^![\Delta_X] = [\Gamma_f]
$$
The last equality follows from
Chow Homology, Lemma \ref{chow-lemma-lci-gysin-easy}.
In other words, going either way around the diagram we
obtain the element of $\text{Corr}^d(X \times Y, \Spec(k))$
corresponding to the cycle $\Gamma_f \subset X \times Y$.
\end{proof}

\end{document}
